\section{Introduction}

There is a classical diagram of Hopf algebras
\be\label{classical-diagram}
\begin{tikzcd}
\Sym \arrow[d, dash] & \NSym \arrow[l, twoheadrightarrow] \arrow[d, dash]  \arrow[r, hook] & \MPR \arrow[d, dash] \\
\Sym \arrow[r,hook] & \QSym  & \MPR  \arrow[l,twoheadrightarrow]
\end{tikzcd}
\ee
in which  the vertical lines are dualities, the $\hookrightarrow$ arrows are inclusions, and 
the $\twoheadrightarrow$ arrows are their adjoints.
The self-dual object $\Sym$ on the left is the familiar Hopf algebra 
of bounded degree symmetric functions \cite[\S2]{GrinbergReiner}, which has an orthonormal basis given by the  Schur functions $s_\lambda$.
The self-dual object $\MPR$ on the right is the \defn{Malvenuto-Reutenauer Hopf algebra}
of permutations from \cite{AguiarSottile,Malvenuto}.
In the middle, we have the dual pair of \defn{quasisymmetric functions} $\QSym$
and \defn{noncommutative symmetric functions} $\NSym$, as described in \cite[\S5]{GrinbergReiner}.

In \cite{LamPyl}, Lam and Pylyavskyy  study a ``$K$-theoretic'' generalization of \eqref{classical-diagram} 
given by 
\be\label{k-diagram}
\begin{tikzcd}
\MSym \arrow[d, dash] & \MNSym \arrow[l, twoheadrightarrow] \arrow[d, dash]  \arrow[r, hook] & \MMPR \arrow[d, dash] \\
\mSym \arrow[r,hook] & \mQSym  & \mMPR  \arrow[l,twoheadrightarrow]
\end{tikzcd}.
\ee
The objects here are modules over $\ZZ[\beta]$ rather than $\ZZ$, where $\beta$ is a formal parameter. Setting $\beta=0$ turns \eqref{k-diagram} into \eqref{classical-diagram}.
The objects $\mSym$ and $\mQSym$ consist of the symmetric and quasisymmetric functions over $\ZZ[\beta]$
of unbounded degree. 
Their duals $\MSym$ and $\MNSym$ are isomorphic to $\Sym$ and $\NSym$ but with scalars extended to $\ZZ[\beta]$.
The objects $\mMPR$ and $\MMPR$, finally, are two different generalizations of 
the Malvenuto-Reutenauer Hopf algebra $
\MPR$.

Besides the Schur functions $\{s_\lambda\}$, the Hopf algebras $\mSym$ and $\MSym$
have another pair of dual bases provided by the \defn{stable Grothendieck polynomials}
$G^{(\beta)}_\lambda$ and the \defn{dual stable Grothendieck polynomials} $g^{(\beta)}_\lambda$.
These power series are relevant to $K$-theory since the $G^{(\beta)}_\lambda$ functions are symmetric limits of connective $K$-theory classes of structure sheaves of Schubert varieties; see \cite{Buch2002,FominKirillov94}.

The goal of this article is to investigate two shifted analogues of \eqref{k-diagram}.
To motivate this, let us first discuss the shifted versions of  \eqref{classical-diagram}.
On one hand, we have a diagram
\be\label{shifted-diagram}
\begin{tikzcd}
\Sym_P \arrow[d, dash] & \Peak_P \arrow[l, twoheadrightarrow] \arrow[d, dash]  \arrow[r, hook] & \MPR \arrow[d, dash] \\
\Sym_Q \arrow[r,hook] & \QSymQ  & \MPR  \arrow[l,twoheadrightarrow]
\end{tikzcd}
\ee
in which the vertical lines are again dualities, the $\hookrightarrow$ arrows are inclusions, and 
the $\twoheadrightarrow$ arrows are their adjoints.
Here  $\Sym_P$ and $\Sym_Q$ are the subalgebras of $\Sym$
spanned by the \defn{Schur $P$-functions} $P_\lambda$ and \defn{Schur $Q$-functions} $Q_\lambda$,
which are indexed by all \defn{strict} integer partitions $\lambda$.
These subalgebras are dual Hopf algebras relative to the bilinear form with $[ P_\lambda, Q_\mu] = \delta_{\lambda\mu}$,
which is different from the usual Hall inner product on $\Sym$;
see \cite[Appendix A]{Stembridge1997a}.
In the middle column, $\QSymQ$ is the Hopf algebra of \defn{peak quasisymmetric functions}
$\QSymQ$ from \cite{Stembridge1997a} while  $\Peak_P$ is the \defn{peak algebra} from \cite{Nyman,Schocker}.

There is another version of \eqref{shifted-diagram} 
in which the roles of $\Sym_P$ and $\Sym_Q$ are interchanged:
\be\label{shifted-diagram2}
\begin{tikzcd}
\Sym_Q \arrow[d, dash] & \Peak_Q \arrow[l, twoheadrightarrow] \arrow[d, dash]  \arrow[r, hook] & \MPR \arrow[d, dash] \\
\Sym_P \arrow[r,hook] & \QSymP  & \MPR  \arrow[l,twoheadrightarrow]
\end{tikzcd}.
\ee
Here $\QSymP$ is a slightly larger version of $\QSymQ$ (namely,  the intersection of $\QSymQ_\QQ := \QQ \otimes_\ZZ \QSymQ$ with $\QSym$ \cite[\S3]{Stembridge1997a})
while its dual $\Peak_Q$ is a Hopf subalgebra of $\Peak_P$.
The diagrams \eqref{shifted-diagram} and \eqref{shifted-diagram2} coincide if we work over $\QQ$ rather than $\ZZ$.


Work of Ikeda and Naruse  \cite{IkedaNaruse} identifies
$K$-theoretic versions $\bGP_\lambda$ and $\bGQ_\lambda$ of the classical Schur $P$- and $Q$-functions.
Whereas $P_\lambda$ and $Q_\lambda$
are generating functions for \defn{(semistandard) shifted tableaux},
$\bGP_\lambda$ and $\bGQ_\lambda$  are generating functions for 
\defn{(semistandard) shifted set-valued tableaux} \cite[Thm.~9.1]{IkedaNaruse}.
These symmetric functions represent the structure sheaves of Schubert varieties in the connective $K$-theory ring of the maximal isotropic Grassmannians of orthogonal and symplectic types \cite[Cor.~8.1]{IkedaNaruse}.

Later results of Nakagawa and Naruse \cite{NN2018} construct
two additional families of ``dual'' $K$-theoretic Schur $P$- and $Q$-functions $\bgp_\lambda$ and $\bgq_\lambda$.
As we will explain in Section~\ref{main-sect}, these power series are $\ZZ[\beta]$-bases 
for two Hopf subalgebras $\MSymP$ and $\MSymQ$ of $\MSym$, 
whose respective duals
 $\mSymQ$ and $\mSymP$
  are 
the completions of the algebras $\ZZ[\beta]\spanning\{\bGQ_\lambda\}$ and $\ZZ[\beta]\spanning\{\bGP_\lambda\}$.
These four objects fit into the pair of
diagrams
\be\begin{aligned}\label{shk-diagram}
\begin{tikzcd}
\MSymP   \arrow[d, dash] &  \MPeakP   \arrow[l, twoheadrightarrow]   \arrow[d, dash]  \arrow[r, hook] & \MMPR \arrow[d, dash] \\
\mSymQ   \arrow[r, hook] &  \mQSymQ   & \mMPR  \arrow[l, twoheadrightarrow]
\end{tikzcd}
\\[-8pt] \\ 
\begin{tikzcd}
\MSymQ   \arrow[d, dash] &  \MPeakQ   \arrow[l, twoheadrightarrow]   \arrow[d, dash]  \arrow[r, hook] & \MMPR \arrow[d, dash] \\
\mSymP   \arrow[r, hook] &  \mQSymP   & \mMPR  \arrow[l, twoheadrightarrow]
\end{tikzcd}
\end{aligned}\ee
which specialize  to \eqref{shifted-diagram} and \eqref{shifted-diagram2} when $\beta=0$,
and which coincide if the scalar ring $\ZZ[\beta]$ is extended to $\QQ[\beta]$.
These diagrams are the primary subject of this article. 
Our main results, building off related work in \cite{ChiuMarberg,LamPyl,LM2022,LM2019}, 
will provide 
distinguished bases for all of the objects shown here, 
explicitly identify the pairings that give the dualities indicated by the vertical lines,
and  
describe the remaining inclusions and their adjoint surjections.

We can summarize our new theorems and outline the rest of this paper as follows.
One way to motivate \eqref{k-diagram}  is through the perspective of \defn{combinatorial Hopf algebras}
as defined in \cite{ABS}. However, some care must be taken to make this rigorous 
as the objects in the bottom row of \eqref{k-diagram} are certain completions of Hopf algebras
rather than actual Hopf algebras.
Section~\ref{prelim-sect} provides a brief survey of the technical background needed to address these issues.



  Section~\ref{k-sect} reviews the construction of the objects and morphisms
in \eqref{k-diagram}. What we present is a very mild generalization of what is studied
by Lam and Pylyavskyy, and involves a parameter $\beta$ that is implicitly set to $\beta=1$ in \cite{LamPyl}.


 The algebras $\mQSymP\supset \mQSymQ$ are spanned by $K$-theoretic generalizations of 
peak quasisymmetric functions studied previously in \cite{LM2019}.
In Section~\ref{multipeak-sect}
we review the definition of these algebras, and prove that $\mQSymQ$ arises of the image of
a canonical morphism of combinatorial Hopf algebras $\mMPR \to \mQSym$; see Theorem~\ref{peak-tl-thm}.

  In Section~\ref{MPeak-sect} we construct the duals of 
$\mQSymP$ and $\mQSymQ$ as   explicit Hopf subalgebras  $\MPeakQ\subset \MNSym$ and $\MPeakP\subset \MNSym$.
This gives two $K$-theoretic generalizations of the classical peak algebra.
Neither appears to have been considered 
in previous literature.
We also derive (co)product formulas for the distinguished bases of 
 $\MPeakQ$ and $\MPeakP$,
 and we identify the adjoint maps $\mMPR \to \mQSymQ$ and $\mMPR \to \mQSymP$ in \eqref{shk-diagram}.
 
  Sections~\ref{shifted-sym-sect} and \ref{pq-sect2}
 discuss the Hopf algebras $\mSymP\supset \mSymQ$
 and their duals $\MSymQ\subset \MSymP$. We 
 prove an identity relating the pairings  $\MSymQ\times \mSymP \to \ZZ[\beta]$
and $\MSymP\times \mSymQ \to \ZZ[\beta]$ to a surjective morphism $\ttheta : \mQSym \to \mQSymQ$;
see Theorem~\ref{ttheta-thm}.
We also identify the adjoint maps $\MPeakP \to \MSymP$
and $\MPeakQ \to \MSymQ$ in \eqref{shk-diagram}; see Theorem~\ref{last-adj-thm}.
These results rely on conjectures from \cite{LM2019,NN2018} proved in \cite{ChiuMarberg,LM2022}.

  Section~\ref{antipode-sect} discusses antipode formulas for the objects in \eqref{shk-diagram}. Building off recent work in \cite{LM2022}, we show that the respective (finite) $\ZZ[\beta]$-linear
spans of all $\bGP$- and $\bGQ$-functions are   sub-bialgebras of $\mSym$ that are not Hopf algebras; see Proposition~\ref{bialgebras-thm}.


  Finally, Section~\ref{open-sect} provides a survey of related open problems and positivity conjectures.
Computer calculations indicate that the coefficients appearing in many different expansions of the distinguished 
bases for the objects in \eqref{shk-diagram} are always positive. In several special cases, it is an open problem to find combinatorial
interpretations of these numbers.
 
 

 

\section{Preliminaries}\label{prelim-sect}

Throughout, we write $\ZZ$ for the set of integers
and let $[n] := \{i \in \ZZ : 0<i \leq n\}$ for $0\leq n \in \ZZ$.
This section presents some basic information about Hopf algebras, their completions,
and quasisymmetric functions.
For more background on each, see \cite{GrinbergReiner,LMW,Marberg2018}.

\subsection{Hopf algebras}

Fix an integral domain $R$ and write $\otimes = \otimes_R$ for the tensor product over this ring.
An \defn{$R$-algebra} is an $R$-module $A$ with $R$-linear
product $\nabla : A\otimes A \to A$
and unit $\iota : R\to A$ 
maps.
Dually, an \defn{$R$-coalgebra} is an $R$-module $A$ with $R$-linear 
coproduct $\Delta : A \to A\otimes A$
and 
counit $\epsilon : A \to R$ maps.
 The (co)product and (co)unit maps must satisfy several associativity axioms;
see \cite[\S1]{GrinbergReiner} for the complete definitions.


An $R$-module $A$ that is both an $R$-algebra and an $R$-coalgebra is an \defn{$R$-bialgebra}
if the coproduct and counit maps are algebra morphisms (equivalently, the product and unit are coalgebra morphisms).

Suppose $A$ is an $R$-bialgebra with structure maps $\nabla$, $\iota$, $\Delta$, and $\epsilon$. 
Let $\End(A)$ denote the set of $R$-linear maps $A \to A$. This set is an $R$-algebra for the 
product 
$
f \ast g := \nabla \circ (f\otimes g) \circ \Delta
$. 
The unit of this \defn{convolution algebra}
is the composition $ \iota\circ \epsilon$ of the unit and counit of $A$.
The bialgebra $A$ is a \defn{Hopf algebra} if  $\id : A \to A$ has a
(necessarily unique) two-sided inverse 
$\antipode : A \to A$ in the convolution algebra $\End(A)$. When it exists, we call $\antipode$ the \defn{antipode} of $A$.

\subsection{Completions}\label{completions-sect}

Many of the objects in the diagrams \eqref{k-diagram} and \eqref{shk-diagram} 
are rings of formal power series of unbounded degree that are ``too large'' to belong to the category of free modules. To formally define algebraic structures on these objects, we need to work in the following setting.





 


 

 

Let $A$ and $B$ be $R$-modules
with an $R$-bilinear form $\langle\cdot,\cdot\rangle : A \times B \to R$.
Assume that $A$ is free and the form
is \defn{nondegenerate} in the sense that
$b\mapsto \langle \cdot,b\rangle$ is a bijection $B\to \Hom_R(A,R)$.
%
Fix a basis $\{a_i\}_{i \in I}$ for $A$.
For each $i \in I$, there exists a unique element $b_i \in B$ with
$\langle a_j, b_i \rangle = \delta_{ij}$ for all $j \in I$,
and we can identify an arbitrary element $b \in B$
with the formal linear combination $\sum_{i \in I} \langle a_i ,b\rangle b_i$.
We refer to $\{b_i\}_{i \in I}$ as a \defn{pseudobasis} for $B$.

We give $R$ the discrete topology.
Then the \defn{linearly compact topology} \cite[\S I.2]{Dieudonne} on $B$ is the coarsest topology 
in which the maps $\langle a_i, \cdot \rangle : B \to R$ are all continuous.
This topology depends on $\langle\cdot,\cdot\rangle$ but not on the choice of basis for $A$,
and is discrete if $A$ has finite rank.




\begin{definition}
A \defn{linearly compact $R$-module} is an $R$-module $B$ equipped with a nondegenerate
bilinear form $A \times B \to R$ for some free $R$-module $A$, given the linearly compact topology;
in this case  $B$ is the \defn{dual} of $A$.
Morphisms between such modules are continuous $R$-linear maps.
\end{definition}

 
 

We will often abbreviate by writing ``LC-'' in place of ``linearly compact.''
Suppose $A$ is a free $R$-module with basis $S$.
We refer to the
$R$-module $B$ of arbitrary $R$-linear combinations of $S$,
 equipped with the nondegenerate bilinear form
$ A \times B \to R$ making $S$ orthonormal,
as the 
\defn{completion} of $A$ with respect to $S$.
This linearly compact $R$-module has $S$ as a pseudobasis.


Let $B$ and $B'$ be linearly compact $R$-modules dual to free $R$-modules $A$ and $A'$.
We reuse $\langle\cdot,\cdot\rangle$ for both of the 
associated nondegenerate forms. 
Every linear map $\phi : A' \to A$ has a unique adjoint $\psi : B\to B'$ such that 
$\langle \phi(a), b\rangle = \langle a,\psi(b)\rangle$ for all $a \in A'$ and $b \in B$.
A linear map $B \to B'$ is \defn{continuous} if and only if it arises as the adjoint 
of some linear map $A' \to A$.




\begin{definition}
The \defn{completed tensor product} of $B$ and $B'$ 
is the $R$-module 
\[B \htimes B' := \Hom_R(A\otimes A',R),\]
given the LC-topology from the tautological pairing 
$(A\otimes A') \times \Hom_R(A\otimes A',R) \to R$.
\end{definition}

If $\{b_i\}_{i \in I}$ and $\{b'_j \}_{j \in J}$ are pseudobases for $B$ and $B'$,
then we can realize $B\htimes B'$ concretely as the linearly compact $R$-module
with the set of tensors $\{ b_i \otimes b_j'\}_{(i,j) \in I \times J}$ as a pseudobasis.



 \begin{example}\label{lc-ex0}
Let $A = R[x]$ and $B=R\llbracket x \rrbracket$.
Define $\langle\cdot,\cdot\rangle : A \times B \to R$ to be the nondegenerate $R$-bilinear form 
\[\left\langle \sum_{n\geq 0} a_n x^n, \sum_{n \geq 0} b_n x^n\right\rangle := \sum_{n\geq 0} a_n b_n.\]
Then the set $\{x^n\}_{n\geq 0}$ is a basis for $A$ and a pseudobasis for $B$,
and we have \[R\llbracket x \rrbracket \otimes R\llbracket y \rrbracket \neq R\llbracket x \rrbracket\htimes R\llbracket y \rrbracket \cong R\llbracket x,y \rrbracket.\]
\end{example}




\begin{definition}
Suppose $\nabla : B \htimes B \to B$ and $\iota : B \to R$ are continuous linear maps
which are the adjoints of linear maps $\epsilon : R \to A$ and $\Delta : A \to A \otimes A$.
We say that $(B,\nabla,\iota)$ is an \defn{LC-algebra}
if $(A,\Delta,\epsilon)$ is an $R$-coalgebra.
Similarly, we say that 
$\Delta : B \to B\htimes B$ and $\epsilon : B \to R$ 
make $B$ into an \defn{LC-coalgebra}
if $\Delta$ and $\epsilon$ are the adjoints of the product and unit 
maps of an $R$-algebra on $A$.
 
\end{definition}

We define
\defn{LC-bialgebras} and \defn{LC-Hopf algebras} analogously.
If $B$ is an LC-Hopf algebra then its \defn{antipode} is the adjoint of the antipode of 
the Hopf algebra $A$.
In each case we say that the (co, bi, Hopf) algebra structures on $A$ and $B$ are duals of each other.

More generally, a Hopf algebra $H$ is \defn{dual} to an LC-Hopf algebra $\hat H$ 
via some nondegenerate bilinear form $\langle\cdot,\cdot\rangle : H \times \hat H \to R$ that is continuous in the second coordinate
if one always has $\langle \iota(a), b\rangle = a\cdot \epsilon(b)$, $\langle a, \iota(b)\rangle = \epsilon(a)\cdot b$,
$\langle \nabla(a \otimes b), c\rangle = \langle a\otimes b,\Delta(c)\rangle$,
and 
$\langle a, \nabla(b \otimes c)\rangle = \langle \Delta(a), b\otimes c \rangle$.


 \begin{example} 
Again let $A = R[x]$ and $B=R\llbracket x \rrbracket$. 
Write $\nabla : B \htimes B \to B$ for the usual product map,
define $\iota : R \to B$ to be the natural inclusion, 
and
let $\epsilon : B \to R$ be the ring homomorphism that sets $x=0$.
For each   $\beta \in R$, there is a continuous algebra homomorphism 
$\Delta_\beta : B \to B \htimes B$ with $\Delta_\beta(x) = x\otimes 1 + 1\otimes x+ \beta x \otimes x$,
and $B$ is an LC-bialgebra relative to $\nabla$, $\iota$, $\Delta_\beta$, and $\epsilon$.

There is a unique bialgebra structure on $A$ that is dual to the one on $B$ via the form in Example~\ref{lc-ex0}.
This structure has unit, counit, and coproduct given by appropriate restrictions of $\iota$, $\epsilon$, and $\Delta_0 := \Delta_\beta|_{\beta=0}$,
while its product  has a more complicated formula; see \cite[Ex. 2.4]{LM2019}.
One can show that the dual bialgebras $A$ and $B$ are dual Hopf algebras:
the antipode of $A$ is the linear map with $\antipode_A(x^n)  = (-1)^ n x(x+\beta)^{n-1}$,
and the antipode of $B$ is the continuous linear map with $\antipode_B(x^n)  =\left(\frac{-x}{1+\beta x}\right)^n$.
Notice that we can restrict $\nabla$ and $\Delta_\beta$ to define a second bialgebra structure on $A$, 
but this will not be a Hopf algebra unless $\beta=0$ as $R$ is an integral domain. 
\end{example}



One can  reformulate the above definitions using commutative diagrams; see 
\cite{Marberg2018}.
Linearly compact (co, bi, Hopf) algebras form a category in which morphisms are continuous linear maps 
commuting with (co)products and (co)units.
The completed tensor product of two linearly compact (co, bi, Hopf) algebras is
itself a linearly compact (co, bi, Hopf) algebra.

   



\subsection{Quasisymmetric functions}

Let $\beta$, $x_1$, $x_2$, \dots be commuting indeterminates.
From this point on, most of our modules will be defined over the ring $R=\ZZ[\beta]$,
and we write $\otimes $ and $\htimes$ for the corresponding tensor products.
A power series $f \in \ZZ[\beta]\llbracket x_1,x_2,\dots\rrbracket $ is \defn{quasisymmetric} if
for any choice of $a_1,a_2,\dots,a_k \in \PP$,
the coefficients of $x_1^{a_1}x_2^{a_2}\cdots x_{k}^{a_k}$ and $x_{i_1}^{a_1}x_{i_2}^{a_2}\cdots x_{i_k}^{a_k}$
in $f$ are equal for all $i_1<i_2<\dots<i_k$.

\begin{definition}
Let $\mQSym$ be the $\ZZ[\beta]$-module of
quasisymmetric power series in $\ZZ[\beta]\llbracket x_1,x_2,\dots\rrbracket $.
Let $\QSym$ denote the submodule of power series in $\mQSym$ of bounded degree.
\end{definition}

A \defn{composition} $\alpha$ is a finite sequence of positive integers.
If the parts of $\alpha$ have sum $N \geq 0$, then we write $\alpha\vDash N$ or $|\alpha|=N$.
The sequence $\alpha$ is a \defn{partition} if it is weakly decreasing, which we indicate by writing $\alpha\vdash N$ instead of $\alpha \vDash N$.

The set $\QSym$ is a graded ring that is free as a $\ZZ[\beta]$-module.
One basis is provided by the
 \defn{monomial quasisymmetric functions}, which are
defined for each composition $\alpha = (\alpha_1,\alpha_2,\dots,\alpha_k)$
as the sums
\[
M_\alpha := \sum_{i_1<i_2<\dots<i_k} x_{i_1}^{\alpha_1} x_{i_2}^{\alpha_2}\cdots x_{i_k}^{\alpha_k} \in \QSym
\] with 
 $M_\emptyset := 1$
when $\alpha =\emptyset$ is the empty composition.
We identify $\mQSym$ with the completion of $\QSym$ relative to this basis.




 Write $\Delta : \QSym \to \QSym \otimes \QSym$ for the $\ZZ[\beta]$-linear map with 
\[
\Delta(M_\alpha) = \sum_{\alpha =\alpha'\alpha''} M_{\alpha'} \otimes M_{\alpha''}
\]
for each composition $\alpha$,
where $\alpha'\alpha''$ denotes the concatenation of   $\alpha'$ and $\alpha''$.
Let $\epsilon : \QSym \to \ZZ[\beta]$ be the linear map with $M_\emptyset \mapsto 1$ and $M_\alpha\mapsto 0$
for all nonempty compositions $\alpha$.
This coproduct and counit make the algebra
 $\QSym$ into a (graded, connected) Hopf algebra \cite[\S5.1]{GrinbergReiner}.
These maps extend to continuous linear maps $\mQSym \htimes \mQSym \to\mQSym$
and
$\mQSym \to \mQSym \htimes \mQSym$ which make $\mQSym$ into an LC-Hopf algebra.
For a description of its antipode, see Section~\ref{antipode-sect}.

Suppose $H $
is an LC-bialgebra, defined over $\ZZ[\beta]$,
with product $\nabla$, coproduct $\Delta$, unit $\iota$, and counit $\epsilon$.
Let $\XX(H)$ be the set of continuous algebra morphisms $\zeta : H \to \ZZ[\beta]\llbracket t\rrbracket $
with $\zeta(\cdot )|_{t=0}=\epsilon$.



\begin{definition}
If $H$ is an LC-bialgebra (respectively, LC-Hopf algebra)
and $\zeta \in \XX(H)$, then  $(H,\zeta)$
is a \defn{combinatorial LC-bialgebra} (respectively, \defn{combinatorial LC-Hopf algebra}).
Such pairs form a category in which morphisms $(H,\zeta) \to (H',\zeta')$
are morphisms $\phi : H \to H'$ with $\zeta = \zeta'\circ \phi$.
\end{definition}



We view $\mQSym$ as a combinatorial LC-Hopf algebra with respect to
the \defn{canonical zeta function} $\zetaq : \mQSym \to \ZZ[\beta]\llbracket t\rrbracket $ 
given by $\zetaq(f) = f(t,0,0,\dots)$.
On monomial quasisymmetric functions, we have 
$\zetaq(M_\alpha)   = t^{|\alpha|}$ for $\alpha\in\{\emptyset,(1),(2),(3),\dots\}$ and $\zetaq(M_\alpha)=0$ for all other compositions $\alpha$.

For each integer $k\geq 1$ define $\Delta^{(k)} := (1 \otimes \Delta^{(k-1)})\circ \Delta = ( \Delta^{(k-1)}\otimes 1 )\circ \Delta$ where $\Delta^{(1)}:=\Delta$. 
Given
a map $\zeta \in \XX(H)$ and a nonempty composition $\alpha
=(\alpha_1,\alpha_2,\dots,\alpha_k)$,
write $\zeta_\alpha : H \to \ZZ[\beta]$ for the map
sending $h \in H$ to the coefficient of 
$t^{\alpha_1}\otimes t^{\alpha_2}\otimes \cdots \otimes t^{\alpha_k}$
in $\zeta^{\otimes k} \circ \Delta^{(k-1)}(h) \in \ZZ[\beta]\llbracket t\rrbracket ^{\htimes k}$.
When $\alpha=\emptyset$ is the empty composition,
define $\zeta_\emptyset := \epsilon$.

\begin{theorem}[{\cite[Thm.~2.8]{LM2019}}]
\label{abs-thm}
Suppose $(H,\zeta)$ is a combinatorial LC-bialgebra.
Then 
the map with the formula
$\Phi(h) = \sum_\alpha \zeta_\alpha(h) M_\alpha$
for $h \in H$,
where the sum is over all compositions $\alpha$,
is the  unique morphism 
of combinatorial LC-bialgebras
$\Phi : (H,\zeta) \to (\mQSym,\zetaq)$.
\end{theorem}









 





 
\section{\texorpdfstring{$K$}{K}-theoretic Hopf algebras}\label{k-sect}

We are now prepared to review the construction of  
 the diagram  \eqref{k-diagram} from \cite{LamPyl}.
As noted in the introduction,
we will work with slightly modified versions of the objects discussed in \cite{LamPyl}, involving a formal parameter $\beta$. 
Setting $\beta=1$ recovers Lam and Pylyavskyy's original definitions, but one can
also go in the reverse direction by making appropriate variable substitutions.


\subsection{Small multipermutations}
\label{wqsym-sect}

We start with the object $\mMPR$ in the lower right corner of \eqref{k-diagram},
called the \defn{small multi-Malvenuto-Reutenauer Hopf algebra} in \cite{LamPyl}.

A \defn{word} is a finite sequence  of positive integers.
Let $v=v_1v_2\cdots v_m$ and $w=w_1w_2\cdots w_n$ be words.
When 
\[S = \{s_1   < \dots <s_m\} \subset [m+n]
\quand
T= \{t_1 <\dots <t_n\} = [m+n]\setminus S,\]
define  $\shuffle_{S,T}(v,w):=u_1u_2\cdots u_{m+n}$
where $u_{s_i} := v_i$  and $u_{t_i} := w_i$.
The \defn{shuffle product} of $v$ and $w$ is 
$v \shuffle w := \sum \shuffle_{S,T}(v,w)$
where the sum is over all pairs $(S,T)$ of disjoint sets with $S\sqcup T =[m+n]$ where $|S| = m$ and $|T|=n$.
For example,  we have  $21 \shuffle 11  
= 3 \cdot 2111 + 2 \cdot 1211 + 1121.$

For $k \in \NN$, let
$w\uparrow k = (w_1+k)(w_2+k)\cdots (w_n+k)$.
If $w$ has $m$ distinct letters, then its standardization is the word
$\st(w) = \phi(w_1)\phi(w_2)\cdots \phi(w_n)$,
where $\phi$ is the unique order-preserving bijection $\{w_1,w_2,\dots,w_n\}\to[m]$.
A word $w$ is \defn{packed}
if $\st(w) = w$.

\begin{definition}
Let $\PWords$ denote the set of packed words and define 
$\mWQSym$ to be the linearly compact $\ZZ[\beta]$-module with $\PWords$ as a pseudobasis. 
\end{definition}

Define   $\nabla : \mWQSym \htimes \mWQSym \to \mWQSym
$
and 
$
\Delta : \mWQSym \to \mWQSym \htimes \mWQSym
$
to be the continuous linear maps with
\be\label{packed-prod-eq}
\ba
\nabla(v\otimes w) &:= v\shuffle (w \uparrow \max(v)),
\\
\Delta(w) &:= \sum_{i=0}^n \st(w_1\cdots w_i) \otimes \st(w_{i+1}\cdots w_n),
\ea
\ee
for $v \in \PWords$ and
$w=w_1w_2\cdots w_n \in \PWords$.
Write 
$\iota : \ZZ[\beta] \to \mWQSym$
and
$\epsilon : \mWQSym\to \ZZ[\beta]$
for the linear maps with
$\iota(1) = \emptyset$
and
$\epsilon(w) = \delta_{w\emptyset}$
for $w \in \PWords$.
These maps make
$\mWQSym$ into an LC-Hopf algebra 
 \cite[Prop.~3.11]{Marberg2018},
 called the Hopf algebra of \defn{word quasisymmetric functions}.


A \defn{small multipermutation} is a packed word with no equal adjacent letters.
Let $\Sm_n$ denote the set of such words $w$ with $\max(w) = n$
and define $\Sm_\infty := \bigsqcup_{n \in \NN} \Sm_n$.
%
Write $<_\m$ for the partial order on $\PWords$ whose cover relations are of the form $w_1\cdots w_i\cdots w_n <_\m w_1 \cdots w_{i}w_{i}\cdots w_n $.
Each lower set under $<_\m$ contains a unique minimal element, which belongs to $\Sm_\infty$.

\begin{definition}
Given  $v \in \Sm_\infty$, let 
$ [v]^{(\beta)}_\m := \sum_{ v \leq_\m w} \beta^{\ell(w)-\ell(v)} w \in \mWQSym$
where the sum is over packed words $w \in \PWords$.
Define $\mMPR$ to be the 
linearly compact $\ZZ[\beta]$-submodule of $\mWQSym$ with  $\left\{[v]^{(\beta)}_\m : v \in \Sm_\infty\right\}$ as a pseudobasis.
\end{definition}



As $\mMPR$ is an LC sub-bialgebra of $\mWQSym$,
which is graded and connected, Takeuchi's formula \cite[Prop.~1.4.22]{GrinbergReiner} implies that its antipode preserves $\mMPR$. This observation lets us
recover the following statement, which is equivalent to \cite[Thms.~4.2 and 7.12]{LamPyl}. 

\begin{theorem}[{\cite{LamPyl}}]
\label{mmpr-thm}
The submodule $\mMPR$ is an LC-Hopf subalgebra of $\mWQSym$.
\end{theorem}



\cite[Thm.~4.2]{LamPyl}
constructs an LC-Hopf algebra over $\ZZ$ with $\Sm_\infty$ as a pseudobasis;
this object is isomorphic to the $\ZZ$-submodule of $\mMPR$ 
with $\left\{\beta^{\ell(v)} [v]^{(\beta)}_\m : v \in \Sm_\infty\right\}$ as a pseudobasis.


 
\subsection{Big multipermutations}

Next, we review the construction of the  \defn{big multi-Malvenuto-Reutenauer Hopf algebra}
from \cite{LamPyl},
which gives the dual object $\MMPR$ in the top right corner of \eqref{k-diagram}.

A \defn{set composition} 
is a sequence of pairwise disjoint nonempty sets $B=B_1B_2\cdots B_m$
with $\bigsqcup_{i \in[m]} B_i = [n]$ for some $n \in \NN$;
in this case we define
 $\ell(B) := m$ and 
 $|B| := n$.

\begin{definition}
Let $\SetComp$ be the set of all set compositions
and define \[\SetComp_n = \{ B \in \SetComp : |B| = n\}\] for $n \in \NN$.
Let $\MWQSym$ be the free $\ZZ[\beta]$-module with $\SetComp$ as a basis.
\end{definition}
%
There is a Hopf algebra structure on $\MWQSym$. 
Given $B=B_1B_2\cdots B_m \in \SetComp$ and $k \in \NN$,
let $k + B $ be the sequence of sets $(k+B_1)(k+B_2)\cdots (k+B_m)$.
For $S \subset \PP$,
define $B \cap S$ 
by removing any empty sets from $(B_1\cap S)(B_2\cap S)\cdots (B_m\cap S)$.
The product $\nabla : \MWQSym \otimes \MWQSym \to \MWQSym$
is the linear map with \[\nabla(A\otimes B) = 
\sum
_{C \in A\bullet B} 
C\]
where
\[A\bullet B := \left\{ C \in \SetComp_{m+n} : C \cap [m] = A\text{ and }C \cap (m+[n]) = m + B\right\}\]
for all $A \in \SetComp_m$ and $B \in \SetComp_n$.
For example,
the elements of $ \{1\}\{2\} \bullet \{1,2\}$ are
$\{1\}\{2\}\{3,4\}$,
$\{1\}\{2,3,4\}$,
$\{1\}\{3,4\}\{2\}$,
$\{1,3,4\}\{2\}$,
and
$\{3,4\}\{1\}\{2\}$.


If $B = B_1B_2\cdots B_m$ is a sequence of subsets of some totally ordered alphabet
and $n=|\bigcup_i B_i|$,
then we let $\st(B):=\phi(B_1) \phi(B_2)\cdots \phi(B_m)$
where $\phi$ is the   order-preserving bijection
$B_1 \cup B_2 \cup \dots \cup B_m \to [n]$.
%
The coproduct
$
\Delta : \MWQSym \to \MWQSym \otimes \MWQSym
$
is   the linear map with \[
\Delta(A) = \sum_{i=0}^m \st(A_1\cdots A_i) \otimes \st(A_{i+1}\cdots A_m)
\quad\text{for all $A=A_1A_2\cdots A_m \in \SetComp$.}\]
Write 
$\iota : \ZZ[\beta] \to \MWQSym$
and
$\epsilon : \MWQSym\to \ZZ[\beta]$
for the linear maps with
$\iota(1) =  \emptyset$
and
$\epsilon(A) = \delta_{A\emptyset}$ for $A \in \SetComp$.
These maps make  $\MWQSym$ into a graded, connected Hopf algebra \cite[\S2.1]{NovelliThibon}.

Consider the following operations interchanging packed words and set compositions.
First, given $w=w_1w_2\cdots w_n \in \PWords$ with $\max(w)=m$,
define $w^*$ to be the set composition $A_1A_2\cdots A_m \in \SetComp_n$
with $A_i = \{ j \in [n] : w_j = i\}$.
Next, for $A =A_1A_2\cdots A_m \in \SetComp_n$, define $A^*$
to be the packed word $w_1w_2\cdots w_n$ with $w_j = i$ if $j \in A_i$.
Finally define $\langle\cdot,\cdot\rangle : \MWQSym \times \mWQSym \to \ZZ[\beta]$ to be the unique bilinear
form,
continuous in the second coordinate, with
\be\label{form1-eq}
\langle A, w \rangle = \delta_{A,w^*}
\qquad\text{for all }A \in \SetComp\text{ and } w\in \PWords.
\ee
 This form is nondegenerate since  $w\mapsto w^*$ and $A \mapsto A^*$
are inverse bijections
$ \PWords \leftrightarrow \SetComp$.
One can also check directly that the relevant products and coproducts
 are compatible 
 in the sense of Section~\ref{completions-sect}.
Therefore $\MWQSym$ and $\mWQSym$ are duals
with respect to \eqref{form1-eq}.


 A \defn{big multipermutation} is a set composition whose blocks never contain
 consecutive integers. 
Let $\SM_n$ be the set of big multipermutations $A$
with $|A| = n$,
and define $\SM_\infty :=\bigsqcup_{n\geq 0} \SM_n$.
The operations $w\mapsto w^*$ and $A \mapsto A^*$
restrict to inverse bijections
$ \Sm_n \leftrightarrow \SM_n$. 

Write $\equivM$ for the partial order 
on set compositions whose cover relations have the form $A \equivM B$ where 
$B$ has a block containing $i$ and $i+1$
and $A = \st(B \cap \{1,\dots,i,i+2,\dots,n\})$.
Each lower set under $\equivM$ contains a unique minimal element, which is a big multipermutation.

\begin{definition}
Let \[
\cJM:=\ZZ[\beta]\spanning\left\{\beta^{|B|-|A|} A- B : A,B\in \SetComp\text{ with }A \equivM B\right\}.\]
Denote the quotient module by $\MMPR:=\MWQSym / \cJM$
and set
\[[A]^{(\beta)}_\M := A + \cJM \in \MMPR\]
for $A \in \SetComp$.
\end{definition}

The $\ZZ[\beta]$-module $\MMPR$ is free  with 
 basis  $\left\{[A]^{(\beta)}_\M: A \in \SM_\infty\right\}$.
One can check that $\cJM$ is the orthogonal complement of $\mMPR$,
which
implies the following results from \cite[\S7.2 and \S7.4]{LamPyl}:

\begin{theorem}[{\cite{LamPyl}}]
\label{form2-thm}
The submodule $\cJM$ is a Hopf ideal of $\MWQSym$,
so $\MMPR$ is a quotient Hopf algebra.
This Hopf algebra is dual to $\mMPR$ via 
the bilinear form
$\langle\cdot,\cdot\rangle : \MMPR \times \mMPR\to \ZZ[\beta]$, 
continuous in the second coordinate, with
\[
\langle [A]^{(\beta)}_\M, [w]^{(\beta)}_\m \rangle = \delta_{A,w^*}
\]
for $A \in \SM_\infty$ and $w\in \Sm_\infty.$
\end{theorem}




The Hopf algebra $\MMPR$ is a very minor generalization of the Hopf algebra constructed in \cite[Thm.~7.5]{LamPyl},
which can be realized inside $\MMPR$ as the $\ZZ$-submodule 
spanned by $  \beta^{|A|} A$ for $A \in \SM_\infty$.




 


 

\subsection{Multifundamental quasisymmetric functions}


Here we review the construction   of an alternate pseudobasis for $\mQSym$,
which arises from viewing $\mMPR$ as a combinatorial LC-Hopf algebra.
The ideas in this section originate in \cite[\S5]{LamPyl}, but we follow the slightly different notational
conventions from \cite[\S3]{LM2019}.
 
For a composition $\alpha=(\alpha_1,\alpha_2,\dots,\alpha_k)\vDash N$
let 
\[
I(\alpha) := \{\alpha_1, \alpha_1+\alpha_2,\dots, \alpha_1+\alpha_2 + \dots + \alpha_{k-1}\}
.\]
Define $\PSet$ to be the set of nonempty, finite subsets of $\PP=\{1,2,3,\dots\}$. 
Given $S,T \in \PSet$, we write $S\preceq T$ if $\max(S) \leq \min(T)$ and $S \prec T$ if $\max(S) < \min(T)$.
\begin{definition}
\label{lbeta-def}
The \defn{multifundamental quasisymmetric function} of $\alpha\vDash N$ is  
\[
 L^{(\beta)}_\alpha   := \sum_S
\beta^{|S| - N} x^S
\in \mQSym
\] where 
the sum is over $N$-tuples $S=(S_1 \preceq S_2 \preceq \dots \preceq S_N)$ with $S_i \in \PSet$
and $S_i \prec S_{i+1}$ if $i \in I(\alpha)$, and where
$|S| := \sum_{i=1}^N |S_i|$ and $x^S := \prod_{i=1}^N \prod_{j \in S_i} x_{j}$.
\end{definition}

The  quasisymmetric functions $ L^{(\beta)}_\alpha$ are another pseudobasis for $\mQSym$ \cite[\S3.3]{LM2019}.

\begin{remark}\label{beta-rmk}
Setting $\beta=0$ transforms $L^{(\beta)}_\alpha$ to
the \defn{fundamental quasisymmetric functions} $L_\alpha := L^{(0)}_\alpha$ \cite[Def.~3.3.4]{LMW}.
Setting $\beta=1$  turns $L^{(\beta)}_\alpha$
into the quasisymmetric functions denoted $\tilde L_\alpha$ in \cite[\S5.3]{LamPyl}.
One recovers $L^{(\beta)}_\alpha$ from $\tilde L_\alpha$ via 
the identity $ L_\alpha^{(\beta)}= \beta^{-|\alpha|}\tilde L_\alpha(\beta x_1,\beta x_2,\dots)$,
which lets one rewrite many formulas in \cite{LamPyl} in terms of $ L_\alpha^{(\beta)}$.
For example, one can obtain explicit expressions
for the product $L^{(\beta)}_{\alpha'} L^{(\beta)}_{\alpha''}$ and coproduct $\Delta(L^{(\beta)}_{\alpha})$
in this way from  \cite[Props. 5.9 and 5.10]{LamPyl}.
\end{remark}

Write $\zeta_{<}$ for the continuous linear map $\mWQSym \to \ZZ[\beta]\llbracket t\rrbracket $
sending strictly increasing packed words $w$ to $t^{\ell(w)}$ and all other packed words to zero.
Then $\zeta_<$ is an algebra morphism with
$\zeta_{<}([w]^{(\beta)}_\m) = \zeta_{<}(w)$ for
all $w \in \Sm_\infty$.
The \defn{descent set} of a word $w=w_1w_2\cdots w_n$ is given by $\Des(w) := \{ i \in [n-1] : w_i > w_{i+1}\}$.
We write $\cC(w)$ 
for the composition 
of $\ell(w)$
with $I(\cC(w)) = \Des(w)$.
The as yet unmotivated definition of  $L^{(\beta)}_\alpha$ is algebraically natural in view of the following:

\begin{theorem}%[{\cite{LamPyl}}]
\label{tl-thm}
The continuous linear map
with $[w]^{(\beta)}_\m \mapsto L^{(\beta)}_{\cC(w)}$ for all $w \in \Sm_\infty$
is
the unique morphism of combinatorial LC-Hopf algebras
$(\mMPR,\zeta_{<}) \to (\mQSym,\zetaq)$.
\end{theorem}

\begin{proof} The claim that this map is a morphism of LC-bialgebras (and therefore also of LC-Hopf algebras) is equivalent to \cite[Thm.~5.11]{LamPyl}.
Choose $w \in \Sm_\infty$ and set $\alpha := \cC(w)$ and $N := |\alpha| = \ell(w)$.
In view of Theorem~\ref{abs-thm}, we just need
to check that $\zeta_{<}([w]^{(\beta)}_\m) = \zetaq(L^{(\beta)}_{\alpha})$.
As $\zetaq$ sends $x_1\mapsto t$ and $x_i\mapsto 0$ for $i>0$,
we either have $\zetaq(L^{(\beta)}_{\alpha}) = t^N$ if the $N$-tuple $S = (\{1\} \preceq \{1\} \preceq \dots \preceq \{1\})$
satisfies the conditions in Definition~\ref{lbeta-def}, or else  $\zetaq(L^{(\beta)}_{\alpha}) =0$.
This means that $\zetaq(L^{(\beta)}_{\alpha}) = t^N = \zeta_{<}([w]^{(\beta)}_\m)$  
if $I(\alpha) = \Des(w)$ is empty and otherwise
$\zetaq(L^{(\beta)}_{\alpha}) = 0 = \zeta_{<}([w]^{(\beta)}_\m)$ as needed.
\end{proof}

 

\subsection{Noncommutative symmetric functions}

We now review the construction from \cite{LamPyl} of the  \defn{multi-noncommutative symmetric functions} $\MNSym$
 in the top row of \eqref{k-diagram}.
%
The \defn{descent set} of a set composition $A = A_1 A_2\cdots A_m \in \SetComp_n$
is 
\[\Des(A) := \{ i  \in [n-1]: i+1 \in A_j\text{ and }i \in A_k\text{ for any indices }j<k\}.\]
One has $\Des(A) = \Des(A^*)$.

\begin{definition}
For a composition $\alpha \vDash n$,
define
\[
\tR_\alpha := \sum_{\Des(A) =I(\alpha) } [A]^{(\beta)}_\M \in \MMPR
\]
where the sum is over big multipermutations $A \in \SM_n$.
These sums are linearly independent, and we
 define $\MNSym$ to be the free $\ZZ[\beta]$-submodule of $\MMPR$
with $\left\{\tR_\alpha : \alpha \text{ is a composition} \right\}$ as a basis.
\end{definition}

 

Recall that we have a form 
$\langle\cdot,\cdot\rangle : \MMPR \times \mMPR\to \ZZ[\beta]$
from Theorem~\ref{form2-thm}.
Evidently if $\alpha$ is a composition and $w \in \Sm_\infty$ has $\gamma = \cC(w)$  
then  $\langle \tR_{\alpha},[w]^{(\beta)}_\m\rangle= \delta_{\alpha\gamma}$.
We reuse the symbol $\langle \cdot,\cdot\rangle $ 
to denote the bilinear form $\langle \cdot,\cdot\rangle : \MNSym \times \mQSym \to \ZZ[\beta]$, continuous in the second coordinate,
 with
$ \langle \tR_{\alpha}, L^{(\beta)}_{\gamma} \rangle = \delta_{\alpha\gamma}$ for all $\alpha$ and $\gamma$.
The following is equivalent to \cite[Thm.~8.4]{LamPyl}:


\begin{theorem}[{\cite{LamPyl}}]
\label{form3-thm}
The module $\MNSym$ is a Hopf subalgebra of $\MMPR$. This subalgebra is the Hopf algebra 
dual to $\mQSym$ via $\langle \cdot,\cdot\rangle $, and the  map $\mMPR \to \mQSym$ 
with 
 $[w]^{(\beta)}_\m \mapsto L^{(\beta)}_{\cC(w)}$ for all $w \in \Sm_\infty$
from Theorem~\ref{tl-thm}
is the morphism adjoint to 
the inclusion
$\MNSym \hookrightarrow \MMPR$.
\end{theorem}


The elements $\tR_{n}:=\tR_{(n)}$ for $n \in \PP$
 freely generate $\MNSym$ as an algebra \cite[Prop.~8.3]{LamPyl},
and one can view $\MNSym$ as a graded connected Hopf algebra in which $\tR_{n}$ has degree $n$.
 In fact, $\MNSym$ is
isomorphic to the usual  \defn{noncommutative symmetric functions} $\NSym$, just defined with scalar ring $\ZZ[\beta]$, by
 \cite[Prop.~8.5]{LamPyl}.

 

\subsection{Symmetric functions}\label{sym-sect}

 A \defn{symmetric function} in $\ZZ[\beta]\llbracket x_1,x_2,\dots\rrbracket $ is a power series 
 that is invariant under permutations of the $x_i$ variables.
The first column of 
 \eqref{k-diagram} contains these familiar  power series:
  
 \begin{definition}
Define $\MSym$ to be the Hopf subalgebra of symmetric functions of bounded degree in $\QSym$.
Let $\mSym$ be the LC-Hopf subalgebra of all symmetric functions in $\mQSym$.
\end{definition}

Let 
$\{s_\lambda\}$ denote the basis of
Schur functions for $\MSym$, indexed by partitions $\lambda$.
It is well-known that $\MSym$ and $\mSym$ are dual Hopf algebras with respect to the nondegenerate bilinear form $\langle\cdot,\cdot\rangle : \MSym \times \mSym \to \ZZ[\beta]$,
 continuous in the second coordinate, that has
$
\langle s_\lambda,s_\mu \rangle = \delta_{\lambda\mu}$.


Lam and Pylyavskyy \cite[Thm.~9.15]{LamPyl} show that $\mSym$ and $\MSym$ have another pair of dual bases given respectively by
 the \defn{stable Grothendieck polynomials} $\{G^{(\beta)}_\lambda\}$
and the \defn{dual stable Grothendieck polynomials} $\{g^{(\beta)}_\lambda\}$,
which satisfy $\langle g^{(\beta)}_\lambda, G^{(\beta)}_\mu\rangle =\delta_{\lambda\mu}$.
 These ``polynomials'' are symmetric generating functions for \defn{semistandard set-valued tableaux}
 and \defn{reverse plane partitions} of shape $\lambda$, respectively.
 
 Since in this article we will never need to work with $G^{(\beta)}_\lambda$ and $g^{(\beta)}_\lambda$ directly,
 we omit their definitions.
 If one does require precise definitions that follow our notational conventions, one 
should adopt the formulas in \cite[Thm.~4.6]{Yeliussizov2019} with $\beta$ replaced by $-\beta$.


 

\begin{theorem}
\label{sym-dual-thm}
The map $\MNSym \to \MSym$ adjoint to the inclusion
$\mSym \hookrightarrow \mQSym$
relative to the forms 
$\langle\cdot,\cdot\rangle$
 is the algebra morphism
with $\tR_{n} \mapsto g^{(\beta)}_{n} := g^{(\beta)}_{(n)}$ for all $n \in \PP$.
\end{theorem}

\begin{proof}
The elements $\tR_{n}$
 freely generate $\MNSym$ \cite[Prop.~8.3]{LamPyl},
 so there is a unique algebra morphism $\MNSym \to \MSym$
 with $\tR_n \mapsto g^{(\beta)}_{n} $ for $n \in \PP$. 
To show that this is the adjoint to $\mSym \hookrightarrow \mQSym$,
it suffices to check that $\langle \tR_{n}, G^{(\beta)}_\lambda \rangle= \delta_{(n),\lambda}$ for all partitions $\lambda$,
as we already know this is the value of $\langle g^{(\beta)}_{n}, G^{(\beta)}_\lambda \rangle $.
As we have $ \langle \tR_{\alpha}, L^{(\beta)}_{\gamma} \rangle = \delta_{\alpha\gamma}$, the desired identity can be deduced from \cite[Eq.\ (3.10)]{LM2019},
which gives the expansion of $G^{(\beta)}_\lambda$ into $L^{(\beta)}_\gamma$'s.
\end{proof}


 \cite[Thm.~9.13]{LamPyl} computes the 
 image of $\tR_\alpha$ 
 under the adjoint map $\MNSym \to \MSym$;
 this  turns out to be a dual stable Grothendieck polynomial indexed by a specific skew ribbon shape.


 





 




 




 

\section{Shifted \texorpdfstring{$K$}{K}-theoretic Hopf algebras}\label{main-sect}

We now turn to the shifted analogues of the diagram \eqref{k-diagram} provided in \eqref{shk-diagram}.
%
We start by 
describing two shifted analogues of $\mQSym$
in Section~\ref{multipeak-sect}.
In Section~\ref{MPeak-sect} we investigate the duals of these LC-Hopf algebras, which provide
$K$-theoretic analogues of the \defn{peak algebra} studied in \cite{Nyman,Schocker}.
Sections~\ref{shifted-sym-sect} and \ref{pq-sect2} 
give an overview of the four (LC-)Hopf algebras of symmetric functions
on the left sides of the two diagrams in.\eqref{shk-diagram}.
In Sections~\ref{antipode-sect} 
we derive several antipode formulas, and then in
Section~\ref{open-sect} we conclude with a survey of open problems.



\subsection{Multipeak quasisymmetric functions}\label{multipeak-sect}

Our first task is to define the shifted analogues of $\mQSym$, which are displayed as
$\mQSymQ$ and $\mQSymP$ in \eqref{shk-diagram}.
This material is partly review from \cite{LM2019}.

For $i \in \ZZ$ let $i' := i - \frac{1}{2}$ so that $\frac{1}{2} \ZZ = \{ \dots < 0' < 0 < 1' < 1 < \dots\}$.
Define $\MSet$ to be the set of finite, nonempty subsets of $\{1'<1<2'<2< \dots\}.$
We again write $S\prec T$ if $\max(S) < \min(T)$ and $S \preceq T$ if $\max(S) \leq \min(T)$
for $S,T \in \MSet$. 
Let 
$x^S := \prod_{i=1}^N \prod_{j \in S_i} x_{\lceil j\rceil }$ and $ |S| := \sum_{i=1}^N |S_i|$
for any sequence $S=(S_1,S_2,\dots,S_N)$ with $S_i \in \MSet$.
A \defn{peak composition} is a composition $\alpha$ 
with $\alpha_i \geq 2$ for $1\leq i < \ell(\alpha)$.
Recall that $I(\alpha) = \{\alpha_1, \alpha_1+\alpha_2, \dots\} \setminus \{|\alpha|\}$. 

\begin{definition}\label{peak-def}
Suppose $\alpha \vDash N$ is a peak composition. Define
\[
K^{(\beta)}_\alpha   := \sum_{
S
}
\beta^{|S|-N} x^S
\]
where the sum is
 over $N$-tuples $S=(S_1 \preceq \dots\preceq S_N)$ of sets $S_i \in \MSet$
with
\be\label{peak-def-eq1}
\ba S_i \cap S_{i+1} &\subset \{1',2',3',\dots\}&&\text{if $i \in I(\alpha)$} \text{ and }
\\
S_i \cap S_{i+1} &\subset \{1,2,3,\dots\}&&\text{if $i \notin I(\alpha)$.}
\ea
\ee
Define \[\bK^{(\beta)}_\alpha   := \sum_{
S
}
\beta^{|S|-N} x^S\] where the sum is over the subset of such $N$-tuples $S$
also satisfying
\be\label{peak-def-eq2}
\ba S_{i+1} &\subset \{1,2,3,\dots\}&&\text{if $i \in \{0\} \sqcup I(\alpha)$.}\ea
\ee
Let $\mQSymQ $ and $\mQSymP$, respectively,
be the  LC-$\ZZ[\beta]$-modules
with  $\{ K^{(\beta)}_\alpha\}$ and $\{ \bK^{(\beta)}_\alpha\}$ (where
$\alpha$ ranges over all peak compositions)
as respective pseudobases.
\end{definition}


The power series $K^{(\beta)}_\alpha$ and $\bK^{(\beta)}_\alpha$ were introduced in \cite{LM2019}
in the context of an ``enriched'' theory of \defn{set-valued $P$-partitions}.
Setting $\beta=0$ transforms $K^{(\beta)}_\alpha$ and $\bK^{(\beta)}_\alpha$ to 
the \defn{peak quasisymmetric functions} defined in \cite[\S3]{Stembridge1997a},
and this implies that   $\{ K^{(\beta)}_\alpha\}$ and $\{ \bK^{(\beta)}_\alpha\}$ are linearly independent.
 However, $K^{(\beta)}_\alpha$ and $\bK^{(\beta)}_\alpha$ are typically not linear combinations 
 (even using rational coefficients and infinitely many terms) of
 the functions $K_\alpha := K^{(0)}_\alpha$
 and 
 $\bK_\alpha := \bK^{(0)}_\alpha$ from \cite[\S3]{Stembridge1997a}.



Both $\mQSymQ$ and $\mQSymP$ are LC-Hopf subalgebras of $\mQSym$,
and
if $\mQSymQ_{\QQ[\beta]}$ is the LC-Hopf algebra defined over $\QQ[\beta]$
 with $\{ K^{(\beta)}_\alpha\}$ as a pseudobasis,
then we have $\mQSymP = \mQSymQ_{\QQ[\beta]} \cap \mQSym\supsetneq \mQSymQ$ \cite[Thm.~4.19]{LM2019}.
More concretely, one has
\be\label{417-eq} K^{(\beta)}_\alpha  = \sum_{\delta \in \{0,1\}^\ell} 2^{\ell-|\delta|} \beta^{|\delta|} \oK^{(\beta)}_{\alpha+\delta}
\quand \oK^{(\beta)}_\alpha  =\sum_{\delta \in (\NN)^\ell} 2^{-\ell-|\delta|}(-\beta)^{|\delta|} K^{(\beta)}_{\alpha+\delta}\ee
for any peak composition $\alpha $ with $\ell=\ell(\alpha)$ parts, where $|\delta| := \sum_{i=1}^\ell \delta_i$  \cite[Cor.~4.17]{LM2019}.
When $\beta=0$, the Hopf algebras $\mQSymQ$ and $\mQSymP$ reduce to (the completions of) the ones denoted
$\mathbf{\Pi}$ and $ \bar{ \mathbf{\Pi}}$ in 
 \cite{BMSW,Stembridge1997a},
 which have been further studied in a number of places (see e.g. \cite{BilleraHsiaoWill,Hsiao,HsiaoPetersen,Li2016,Petersen}).


Recall the definition of $\zeta_{<}:\mWQSym \to \ZZ[\beta]\llbracket t\rrbracket $ 
and write 
$\zeta_{>}:\mWQSym \to \ZZ[\beta]\llbracket t\rrbracket $
for the continuous linear map 
whose value at $w=w_1\cdots w_n \in \PWords$ is
\be
\zeta_>(w_1w_2\cdots w_n) := \zeta_<(w_n \cdots w_2w_1)
= 
\begin{cases} t^n &\text{if }w_1>w_2>\dots> w_m \\ 0&\text{otherwise}.\end{cases}
\ee
This is an algebra morphism
with $\zeta_{>}([w]^{(\beta)}_\m) = \zeta_{>}(w)$ for $w \in \Sm_\infty$.
By Theorem~\ref{abs-thm} there is a unique morphism of combinatorial LC-Hopf algebras
$(\mMPR,\zeta_{>}) \to (\mQSym,\zetaq)$.
Although this map is different from the one in Theorem~\ref{tl-thm},
it also sends $\{ [w]^{(\beta)}_\m : w \in \Sm_\infty\}$ 
to the pseudobasis of multifundamental quasisymmetric functions $\{ L^{(\beta)}_\alpha\}$; see the proof of \cite[Prop.~6.3]{LM2019}.

To construct something new, we consider the \defn{convolution} of the maps $\zeta_{>}$ and $\zeta_{<}$ defined by
the formula
$ \zeta_{>|<} := \nabla_{\ZZ[\beta]} \circ (\zeta_> \htimes \zeta_<)\circ \Delta : \mWQSym \to \ZZ\llbracket t\rrbracket . $
This is a continuous algebra morphism $ \mWQSym \to \ZZ\llbracket t\rrbracket $.
If $w=w_1w_2\cdots w_n \in \PWords$ then 
\be \zeta_{>|<}(w)=
\begin{cases}
2 t^n &\text{if $w_1>\dots > w_i <\dots < w_n$ for some $i \in [n]$} \\
t^n &\text{if $w_1>\dots > w_i = w_{i+1} < \dots < w_n$ for some $i \in [n-1]$} \\
1 & \text{if }n=0\\
0&\text{otherwise}.\end{cases}
\ee
It follows that if $w=w_1w_2\cdots w_n \in \Sm_\infty$ then
\be \zeta_{>|<}([w]^{(\beta)}_\m) = \begin{cases}
t^n(2+\beta t) &\text{if $w_1>\dots > w_i < \dots < w_n$ for some $i \in [n]$} \\
1 &\text{if } n=0\\
0&\text{otherwise}.
\end{cases} 
\ee

For any composition $\alpha\vDash n$, let $\Lambda(\alpha)$ be the 
unique peak composition of $n$ satisfying $ I(\Lambda(\alpha)) = \{ i \in I(\alpha) : 0 < i-1 \notin I(\alpha)\}.$
For example, one has \[\Lambda((1,2,1,1,1,3,1)) = (3,6,1).\] 
Then define
$\ttheta : \mQSym \to \mQSymQ$
to
be the continuous linear map with 
\be\label{ttheta-eq} 
\ttheta(L^{(\beta)}_\alpha) = K^{(\beta)}_{\Lambda(\alpha)}
\quad\text{for all compositions $\alpha$.}
\ee
By \cite[Cor.~4.22]{LM2019},
this map is a surjective morphism of LC-Hopf algebras.

The \defn{peak set} of a word $w=w_1w_2\cdots w_n$ is 
\[\PeakSet(w) := \{1<i<n : w_{i-1} < w_i > w_{i+1}\}.\]
Let $\cD(w)$ be the unique peak composition $\alpha\vDash\ell(w)$ with $I(\alpha) = \PeakSet(w)$.
If $w \in \Sm_\infty$ then 
\be\label{peakset-eq}
\PeakSet(w) = \{ i \in \Des(w) : 0 < i-1 \notin \Des(w)\}\quad\text{so}\quad 
\Lambda(\cC(w))=   \cD(w),\ee
and we have
$ \zeta_{>|<}([w]^{(\beta)}_\m) \neq 0$ if and only if $\PeakSet(w) = \varnothing$, in which case $\ell(\cD(w))\leq 1$.
The multipeak quasisymmetric functions are motivated algebraically by this analogue of Theorem~\ref{tl-thm}:


\begin{theorem}\label{peak-tl-thm}
The continuous linear map with $[w]^{(\beta)}_\m \mapsto K^{(\beta)}_{\cD(w)}$ for $w \in \Sm_\infty$
is the unique morphism of 
combinatorial LC-Hopf algebras 
$(\mMPR,\zeta_{>|<}) \to (\mQSym, \zetaq)$.
\end{theorem}

\begin{proof}
If $w \in \Sm_\infty$ then 
 $\cD(w) = \Lambda(\cC(w))$ since
 \[
 \ba 
 I( \Lambda(\cC(w))) &= \{ i \in I(\cC(w)) : 0 < i-1 \notin I(\cC(w))\}\\& =  \{ i \in \Des(w) : 0 < i-1 \notin \Des(w)\}= \PeakSet(w)  .\ea\]
Our map $\Psi : \mMPR \to \mQSym$ is thus the composition of 
$\Phi : (\mMPR,\zeta_{<}) \to (\mQSym, \zetaq)$ from Theorem~\ref{tl-thm} and $\ttheta : \mQSym \to \mQSymQ$,
so $\Psi$ is at least a morphism of LC-Hopf algebras. 

It remains to check that $\zeta_{>|<} = \zetaq\circ \Psi$. For this,  it 
suffices to show that if $\alpha\vDash N$ is a peak composition then
  $ \zetaq(K^{(\beta)}_\alpha) = t^{|\alpha|}(2+\beta t)$ if $I(\alpha) = \varnothing$ and $ \zetaq(K^{(\beta)}_\alpha) =0$ otherwise.
Recall that $\zetaq$ corresponds to setting $x_1=t$ and $x_i=0$ for $i>1$.
Thus
$
\zetaq(K^{(\beta)}_\alpha)  = \sum_{S} \beta^{|S|-N} t^{|S|}
$ where the sum is over all weakly increasing $N$-tuples of sets $S= (S_1\preceq S_2 \preceq \cdots \preceq S_N)$ with $\varnothing\neq S_i \subseteq \{1'<1\}$,  $S_i \cap S_{i+1} \subseteq \{1'\}$ for $i \in I(\alpha)$, and $S_i \cap S_{i+1} \subseteq \{1\}$ for $i \notin I(\alpha)$. 

If $I(\alpha) \neq \varnothing$,  then $\alpha_1 \geq 2$ since $\alpha$ is a peak composition, so there are no such tuples $S$ since we must have $S_{i-1} \cap S_i \subset \{1\}$ and $S_i \cap S_{i+1} \subset \{1'\}$ for $i \in I(\alpha)$. Thus $\zetaq(K^{(\beta)}_\alpha) = 0$ if $I(\alpha) \neq \varnothing $ as claimed.
%
On the other hand, if $I(\alpha) = \varnothing$, then we have $S_i \cap S_{i+1} \subseteq \{1\}$ for all $i \in [N-1]$, so there
are only three possibilities for $S$, given by 
$( \{1\}, \{1\},\dots,\{1\})$, $( \{1'\}, \{1\},\dots,\{1\})$, and  $( \{1',1\}, \{1\},\dots,\{1\})$,
so we have $ \zetaq(K^{(\beta)}_\alpha) = t^{|\alpha|}(2+\beta t)$ as needed.
\end{proof}






At this point it is useful to describe the product and coproduct in $\mMPR$ more explicitly.
Recall the definition of $\leq_\m$ from Section~\ref{wqsym-sect}.
Given a word $w$ and a set $S$, define $w\cap S$ to be the subword of $w$ formed by omitting all letters not in $ S$.
Then, as explained in \cite[\S4]{LamPyl}, one has
\be\label{mMPR-prod-eq}
 \nabla([w']^{(\beta)}_\m\otimes [w'']^{(\beta)}_\m) = \sum_w  \beta^{\ell(w) - \ell(w')-\ell(w'')}[w]^{(\beta)}_\m
\quad\text{for $w' \in \Sm_m$ and $w'' \in \Sm_n$,}\ee
where the sum is over all $w \in \Sm_{m+n}$ such that $w' \leq_\m w \cap [m] $ and $w'' \uparrow m \leq_\m  w \cap (m+[n])$.
If
we
fix $w =w_1w_2\cdots w_n \in \Sm_\infty$ and define $\llbracket v\rrbracket ^{(\beta)}_\m := [\st(v)]^{(\beta)}_\m$ for an arbitrary word $v$,
then
\be\label{mMPR-coprod-eq}
\ba \Delta([w]^{(\beta)}_\m) = &\sum_{i=0}^n  \llbracket w_1\cdots w_i\rrbracket ^{(\beta)}_\m  \otimes \llbracket w_{i+1}\cdots w_n\rrbracket ^{(\beta)}_\m
\\&+ \beta \sum_{i=1}^n \llbracket w_1\cdots w_i\rrbracket ^{(\beta)}_\m  \otimes \llbracket w_i\cdots w_n\rrbracket ^{(\beta)}_\m.
\ea\ee
These formulas follow directly from the definitions of $\mWQSym$ and $[w]^{(\beta)}_\m$.
Using Theorem~\ref{form2-thm}, one can translate these identities by duality to product and coproduct formulas
for $\MMPR$; see \cite[\S7.1]{LamPyl}.
On the other hand, invoking Theorem~\ref{peak-tl-thm} leads to the following formulas for $\mQSymQ$:



\begin{proposition}\label{K-prod-prop}
Suppose $\alpha' $ and $\alpha'' $ are peak compositions. Choose any
 $w' ,w'' \in \Sm_\infty$ with $\cD(w') = \alpha'$ and $\cD(w'') =\alpha''$,
 and set $m=\max(\{0\} \cup w')$ and $n=\max(\{0\} \cup w'')$.
  Then
\[
K^{(\beta)}_{\alpha'} K^{(\beta)}_{\alpha''} = \sum_w \beta^{\ell(w) - |\alpha'| - |\alpha''|} K^{(\beta)}_{\cD(w)}
\]
where the sum is over all $w \in \Sm_{m+n}$ with $w' \leq_\m w \cap [m]$ and $w'' \uparrow m \leq_\m w \cap (m + [n])$.
%
If $\alpha$ is a peak composition and $w=w_1w_2\cdots w_\ell \in\Sm_\infty$ has  
$\cD(w) = \alpha$, then
\[
\Delta(K^{(\beta)}_{\alpha}) = \sum_{i=0}^\ell K^{(\beta)}_{\cD(w_1\cdots w_i)} \otimes K^{(\beta)}_{\cD(w_{i+1}\cdots w_\ell)}
+
\beta \sum_{i=1}^\ell K^{(\beta)}_{\cD(w_1\cdots w_i)} \otimes K^{(\beta)}_{\cD(w_{i}\cdots w_\ell)}.
\]
\end{proposition}

\begin{proof}
Apply the morphism in Theorem~\ref{peak-tl-thm} to both sides of  \eqref{mMPR-prod-eq} and \eqref{mMPR-coprod-eq}.
\end{proof}


In principle one can also
 compute products and coproducts in the pseudobasis $\{ \bK^{(\beta)}_\alpha\}$
by combining the formulas in Proposition~\ref{K-prod-prop} with the change-of-basis identities in \eqref{417-eq}.

\begin{example}\label{K-prod-ex}
If $\alpha' = \alpha'' = (1)$ then $K^{(\beta)}_{\alpha'} K^{(\beta)}_{\alpha''}
 $
is the sum $ \sum_{w} \beta^{\ell(w) - 2} K^{(\beta)}_{\cD(w)}$ over all words  $w\in \{ 12, 21, 121, 212, 1212, 2121, 12121,21212, \dots\},$ 
so there is an infinite product expansion 
\[K^{(\beta)}_{(1)} K^{(\beta)}_{(1)} 
=
 2 K^{(\beta)}_{(2)} + \beta K^{(\beta)}_{(2,1)} + \beta K^{(\beta)}_{(3)}+ \beta ^2 K^{(\beta)}_{(2,2)}+ \beta^2 K^{(\beta)}_{(3,1)}+ 
 \beta^3 K^{(\beta)}_{(2,2,1)}%+\beta^3 K^{(\beta)}_{(3,2)}
  + \cdots .
\] However, there is a finite coproduct expansion
\[ \ba
\Delta(K^{(\beta)}_{(2)}) &= \Delta(K^{(\beta)}_{\cD(12)})
\\& =
1 \otimes K^{(\beta)}_{(2)} 
+ K^{(\beta)}_{(1)}\otimes K^{(\beta)}_{(1)}
+  K^{(\beta)}_{(2)}\otimes 1
+
\beta\left(  K^{(\beta)}_{(1)}\otimes K^{(\beta)}_{(2)}
+
  K^{(\beta)}_{(2)}\otimes K^{(\beta)}_{(1)}\right).\ea
 \]
\end{example}





 

 


\subsection{Multipeak noncommutative symmetric functions}\label{MPeak-sect}

We now consider the duals of $\mQSymQ$ and $\mQSymP$.
%
Fix a set composition $A = A_1 A_2\cdots A_m \in \SetComp_n$,
so that the union of the blocks of $A$ is $[n]$.
Recall that $i \in [n-1]$ belongs to $\Des(A)$
if and only if the block of $A$ containing $i$ is after the block containing $i+1$. 

The \defn{peak set} of $A$ is
$\PeakSet(A) := \{1<i<n: i -1 \notin \Des(A)\text{ and } i \in \Des(A) \}.$
If $A $ belongs to $ \SM_n$ (so that none of its blocks contain consecutive integers)
 then one has $i \in \PeakSet(A)$ precisely when $1< i <n$ and the block of $A$ containing $i$ is after the 
blocks containing $i-1$ and $i+1$.
Even when $A \notin \SM_n$, the set $\PeakSet(A)$ is always equal to $ I(\alpha)$ for some peak composition $\alpha \vDash n$.




\begin{definition}
For a peak composition $\alpha \vDash n$,
let
\[\tpeakR_\alpha := \sum_{\substack{A \in \SM_n \\  \PeakSet(A) =I(\alpha)}} [A]^{(\beta)}_\M  \in \MMPR.
\]
Define $\MPeakP$ to be the free $\ZZ[\beta]$-module 
with $\left\{\tpeakR_\alpha : \alpha\text{ is a peak composition}\right\}$ as a basis.
\end{definition}

 

It also holds that $\tpeakR_\alpha =  \sum_{\Lambda(\gamma) =\alpha} \tR_\gamma$
where the sum is over $\gamma \vDash |\alpha|$, so $\MPeakP\subseteq\MNSym$.
Define
$ [\cdot,\cdot ] : \MPeakP \times \mQSymQ \to \ZZ[\beta] $
to be the nondegenerate bilinear form, continuous in the second coordinate, that has 
$[ \tpeakR_\alpha, K^{(\beta)}_\gamma] = \delta_{\alpha\gamma}$
for all peak compositions $\alpha$ and $\gamma$.
Below, let $\langle \cdot,\cdot\rangle : \MNSym \times \mQSym \to \ZZ[\beta]$
be as in Theorem~\ref{form3-thm} 
and recall the definition of $\ttheta $ from \eqref{ttheta-eq}.


 

\begin{lemma}\label{[]-lem}
If $f \in \MPeakP$ and $g \in \mQSym$
then $\left[ f, \ttheta(g)\right] = \langle f,g\rangle$. 
\end{lemma}

\begin{proof}
We may assume that $f = \tpeakR_\alpha$ and $g = L^{(\beta)}_\gamma$
 for a peak composition $\alpha$ and a composition $\gamma$. 
Then the desired identity is clear by comparing the definitions of $\ttheta$ and $\tpeakR_\alpha$.
\end{proof}


\begin{theorem}\label{mpeakp-thm}
The module $\MPeakP$ is a Hopf subalgebra of $\MNSym$
and is the Hopf algebra dual to $\mQSymQ$ via   $[\cdot,\cdot]$.
The continuous linear map 
$\mMPR \to \mQSymQ$ 
from Theorem~\ref{peak-tl-thm}
sending
$[w]^{(\beta)}_\m \mapsto K^{(\beta)}_{\cD(w)}$  for all $w \in \Sm_\infty$
is the morphism adjoint to 
the inclusion
$\MPeakP \hookrightarrow \MMPR$.
\end{theorem}


\begin{proof}
Relative to the form in Theorem~\ref{form2-thm},
the set $\MPeakP$ is the orthogonal complement
of the kernel of the LC-Hopf algebra morphism 
$\mMPR \to \mQSymQ$ described in Theorem~\ref{peak-tl-thm}.
Therefore $\MPeakP$ is a Hopf subalgebra.
Lemma~\ref{[]-lem}, in view of Theorem~\ref{form3-thm}, implies
that the nondegenerate form $[\cdot,\cdot]$ respects the (co)product and (co)unit maps of 
$\MPeakP$ and $\mQSymQ$, so
$\MPeakP$  dual to $\mQSymQ$.
For the last assertion, 
we note that if $\alpha$ is a peak composition and $w \in \Sm_\infty$
then  
\[\langle\tpeakR_\alpha, [w]^{(\beta)}_\m\rangle
=\langle 
  \tpeakR_\alpha, L^{(\beta)}_{\cC(w)} \rangle 
    =[ 
  \tpeakR_\alpha, \ttheta(L^{(\beta)}_{\cC(w)} )]    
   =[ 
  \tpeakR_\alpha, K^{(\beta)}_{\cD(w)} ] 
  \]
  by Theorem~\ref{form3-thm} for the first equality,  
   Lemma~\ref{[]-lem} for the second, and \eqref{peakset-eq} for the third.
\end{proof}


We call $\MPeakP$ the \defn{multi-peak algebra}. 
This is a generalization of the \defn{peak algebra} carefully studied in \cite{Schocker} (see also \cite{AguiarBergeronNyman,AguiarNymanOrellana,BergeronHivertThibon,JingLi,Nyman}),
which coincides with $\MPeakP$ when $\beta=0$.

We can compute a product formula for the $\tpeakR_\alpha$-basis of $\MPeakP$.
Suppose
$\alpha$ and $\gamma$ are nonempty peak compositions of length $m$ and $n$.
Define 
$\alpha \vartriangleleft \gamma:= \alpha\gamma$ and 
\[
\ba
\alpha \vartriangleright \gamma &:=  (\alpha_1,\dots,\alpha_{m-1},\alpha_m+ \gamma_1,\gamma_2,\dots,\gamma_n),
\\
\alpha \circ \gamma &:=  (\alpha_1,\dots,\alpha_{m-1},\alpha_m+ \gamma_1-1,\gamma_2,\dots,\gamma_n).
\ea
\]
Additionally let
\be\ba
\alpha \blacktriangleright \gamma &:= \alpha \vartriangleright (1, \gamma_1-1,\gamma_2,\dots,\gamma_n)\\& =  (\alpha_1,\dots,\alpha_{m-1},\alpha_m+1, \gamma_1-1,\gamma_2,\dots,\gamma_n),
\\
\alpha\bullet \gamma &:=  \alpha \circ (1, \gamma_1-1,\gamma_2,\dots,\gamma_n)  \\&= (\alpha_1,\dots, \alpha_m, \gamma_1-1,\gamma_2,\dots,\gamma_n).
\ea\ee
If $n=1$ then $\alpha \blacktriangleright \gamma$ and $\alpha\bullet \gamma $ could be integer sequences ending in zero,
which we do not consider to be peak compositions.
We define $\tpeakR_\alpha := 0$ if $\alpha$ is not a peak composition.
The following result is a shifted analogue of the product formula \cite[Prop.~8.1]{LamPyl}
for the $\tR_{\alpha}$-basis of $\MNSym$.

\begin{proposition}\label{MPeakP-prod-prop}
Suppose $\alpha$ and $\gamma$ are nonempty peak compositions. Then
\[\tpeakR_{\alpha}  \tpeakR_{\gamma}
=  \tpeakR_{\alpha\blacktriangleright \gamma} + \tpeakR_{\alpha\vartriangleright \gamma} + \tpeakR_{\alpha \vartriangleleft \gamma} 
+ \beta\cdot \tpeakR_{\alpha \circ \gamma} +  \beta\cdot  \tpeakR_{\alpha \bullet \gamma}.
\]
\end{proposition}



  \begin{proof}
  Choose a word $w=w_1w_2\cdots w_n \in \Sm_\infty$.
By Proposition~\ref{K-prod-prop} and Theorem~\ref{mpeakp-thm}   
the expression ${[\tpeakR_{\alpha}  \tpeakR_{\gamma}, K^{(\beta)}_{\cD(w)}]} = 
[\tpeakR_{\alpha} \otimes \tpeakR_{\gamma}, \Delta(K^{(\beta)}_{\cD(w)})] $ expands to \[
\sum_{i=0}^n \delta_{\alpha,\cD(w_1\cdots w_i)} \delta_{\gamma,\cD(w_{i+1}\cdots w_n)}
+
\beta \sum_{i=1}^n \delta_{\alpha, \cD(w_1\cdots w_i)} \delta_{\gamma,\cD(w_{i}\cdots w_n)}.
\]
Now observe that one can have $\alpha=\cD(w_1\cdots w_i) $ and $\gamma=\cD(w_{i+1}\cdots w_n)$ 
precisely when either $i \in \PeakSet(w)$ and $\cD(w) = \alpha\vartriangleleft \gamma$,
$i+1 \in \PeakSet(w)$ and $\cD(w) = \alpha\blacktriangleright \gamma$,
or $\{i,i+1\} \cap \PeakSet(w) = \varnothing$ and $\cD(w) =\alpha\vartriangleright\gamma$.
Similarly, one can have 
$\alpha=\cD(w_1\cdots w_i) $ and $\gamma=\cD(w_{i}\cdots w_n)$
precisely when
 $i \in \PeakSet(w)$ and $\cD(w) = \alpha\bullet \gamma$
 or when
 $i \notin \PeakSet(w)$ and $\cD(w) = \alpha\circ \gamma$.
  \end{proof}

\begin{example}
Here are two examples of Proposition~\ref{MPeakP-prod-prop}.  All five terms appear in  
%\[
%\tpeakR_{(3,2,5,2)}  \tpeakR_{(4,2)} = \tpeakR_{(3,2,5,3,3,2)} + \tpeakR_{(3,2,5,6,2)}  + \tpeakR_{(3,2,5,2,4,2)} +\beta \cdot\tpeakR_{(3,2,5,5,2)} +\beta \cdot\tpeakR_{(3,2,5,2,3,2)}.
%  \]
\[
\tpeakR_{3252}  \tpeakR_{42} = \tpeakR_{325332} + \tpeakR_{32562}  + \tpeakR_{325242} +\beta \cdot\tpeakR_{32552} +\beta \cdot\tpeakR_{325232}.
  \]
On the other hand, only three survive in 
%\[
%\ba \tpeakR_{(3,2,5,1)}  \tpeakR_{(4,2)} &= \tpeakR_{(3,2,5,2,3,2)} + \tpeakR_{(3,2,5,5,2)}  + \tpeakR_{(3,2,5,1,4,2)} +\beta \cdot\tpeakR_{(3,2,5,4,2)} +\beta\cdot \tpeakR_{(3,2,5,1,3,2)}
%\\
%&=
% \tpeakR_{(3,2,5,2,3,2)} + \tpeakR_{(3,2,5,5,2)}  +\beta\cdot \tpeakR_{(3,2,5,4,2)} .
%\ea
%  \]
\[
\ba \tpeakR_{3251}  \tpeakR_{42} &= \tpeakR_{325232} + \tpeakR_{32552}  + \tpeakR_{325142} +\beta \cdot\tpeakR_{32542} +\beta\cdot \tpeakR_{325132}
\\
&=
 \tpeakR_{325232} + \tpeakR_{32552}  +\beta\cdot \tpeakR_{32542} .
\ea
  \]
To save space here, we have removed the parentheses and commas from all subscripts ``$(\alpha_1,\alpha_2,\dots,\alpha_l)$'' and rewritten these as ``$\alpha_1\alpha_2\cdots \alpha_l$''.
   \end{example}
 


Below, we abbreviate by writing $\tpeakR_{n}:=\tpeakR_{(n)}$ for $n \in \PP$ and set $\tpeakR_0 := 1$.

\begin{lemma}\label{peak-free-lem}
If $n$ is a positive integer and $\delta_{[n]} := | \{ n\} \cap \{2,4,6,\dots\}|$ then
\[ \tpeakR_1 \tpeakR_{n-1} \in \delta_{[n]}  \cdot \tpeakR_n +   \sum_{i=2}^{n-1} (-1)^i \tpeakR_i \tpeakR_{n-i}  
+  \ZZ[\beta]\spanning\left\{ \tpeakR_\alpha : |\alpha| <  n\right\}.
\]
\end{lemma}

\begin{proof}
Let $\cI =   \ZZ[\beta]\spanning \{ \tpeakR_\alpha : |\alpha| <  n \}$.
Proposition~\ref{MPeakP-prod-prop}
implies 
$\tpeakR_1 \tpeakR_{n-1}  +\cI = \tpeakR_{n} +\tpeakR_{(2,n-2)} + \cI
$
for $n>1$, where $\tpeakR_{(n,0)} := 0$.
Likewise, if $2\leq i < n$ then
$
\tpeakR_{(i, n-i)} +\cI = \tpeakR_{i}\tpeakR_{n-i}- \tpeakR_{n}-\tpeakR_{(i+1,n-i-1)} + \cI
$.
We obtain the lemma by 
successively expanding the right hand side of the first identity using the second. 
\end{proof}

\begin{proposition}\label{MPeak-free-prop}
The set $\left\{\tpeakR_n: n =1,3,5,\dots\right\}$ freely generates $\MPeakP$ as a $\ZZ[\beta]$-algebra.
\end{proposition}

\begin{proof}
Suppose $\alpha=(\alpha_1,\alpha_2,\dots,\alpha_m)$ is a composition
and let \[\bXi_\alpha := \tpeakR_{\alpha_1}  \tpeakR_{\alpha_2} \cdots  \tpeakR_{\alpha_m}.\]
We say that $\alpha$ is \defn{odd} if $\alpha_i$ is an odd integer for each $i \in [m]$.
It suffices to show that the elements  $\bXi_\alpha$ form a $\ZZ[\beta]$-basis for $\MPeakP$ when $\alpha$ ranges over all odd compositions.
%

First let $<_{\mathrm{revlex}}$ be the partial order on compositions with $\gamma <_{\mathrm{revlex}} \alpha$ 
if $|\gamma| < |\alpha|$ or if $|\gamma| = |\alpha|$ and $\gamma$ exceeds $\alpha$ in lexicographic order.
It follows from Proposition~\ref{MPeakP-prod-prop}
that if $\alpha$ is any peak composition then $\bXi_\alpha \in \tpeakR_\alpha + \ZZ[\beta]\spanning\left\{ \tpeakR_\gamma : \gamma <_{\mathrm{revlex}} \alpha\right\}$.
Thus  the elements  $\bXi_\alpha$ form a basis for $\MPeakP$ at least when $\alpha$ ranges over all peak compositions.

If $\alpha$ is a peak composition then let $\mathsf{odd}(\alpha)$ be the odd composition formed by replacing 
each even part $\alpha_i$ by two consecutive parts $(1,\alpha_i-1)$. 
For example, $\mathsf{odd}((3,6,3,4,2)) = (3,1,5,3,1,3,1,1)$.
Let $<_{\mathrm{lex}}$ be the partial order on compositions with $\gamma <_{\mathrm{lex}} \alpha$ if $|\gamma| < |\alpha|$
or if $|\gamma| = |\alpha|$ and $\gamma$ precedes $\alpha$ lexicographically.
It follows by induction on $\ell(\alpha)$ using Lemma~\ref{peak-free-lem}
that
$\bXi_{\mathsf{odd}(\alpha)} \in \bXi_\alpha + \ZZ[\beta]\spanning\left\{ \bXi_\gamma : \gamma <_{\mathrm{lex}} \alpha\right\}$.
Since $\mathsf{odd}$ is a bijection from peak compositions to odd compositions,
we deduce that $\left\{\bXi_\alpha : \alpha\text{ is an odd composition}\right\}$
 is another $\ZZ[\beta]$-basis for $\MPeakP$ as desired.
\end{proof}
